% TOP_PROBLEM: {"id": 2923, "title": "Kirby Problem 4.47", "category_id": 7, "category": {"id": 7, "name": "topology", "display_name": "Topology", "description": "Properties preserved under continuous deformations.", "slug": "topology", "order_index": 7, "created_at": "2026-07-31T15:26:25.671Z"}, "status": "open"}
% FIRST_POSTED: null
% ENTRY_KIND: statement_only
% Imported from: batch13.tex; UnsolvedMath ID 2923
\documentclass[11pt]{article}
\usepackage[margin=25mm]{geometry}
\usepackage{fontspec}
\setmainfont{Latin Modern Roman}
\usepackage{amsmath,amssymb,mathtools}
\usepackage{unicode-math}
\setmathfont{Latin Modern Math}
\usepackage{xurl,hyperref}
\hypersetup{colorlinks=true,urlcolor=blue}
\setcounter{secnumdepth}{0}
\setlength{\parindent}{0pt}
\setlength{\parskip}{5pt}
\setlength{\emergencystretch}{3em}
\newcommand{\sref}[2]{\hyperlink{source:#1:#2}{[S#2]}}
\newcommand{\cataloguescope}[1]{\paragraph{Recorded target.}#1}
\begin{document}
% UnsolvedMath ID 2923; source catalogue code KP-4.47
\setcounter{section}{2922}
\section{Kirby Problem 4.47}\label{2923}
{\small\sffamily UnsolvedMath ID 2923 · Topology}
\subsection{Definitions and mathematical statement}

\paragraph{Shared notation.}$\mathbb N=\{1,2,\ldots\}$, and $\mathbb Z,\mathbb Q,\mathbb R,\mathbb C$ denote the integers, rationals, reals and complex numbers. Any other conventions are specified locally.


Let $L=L_1\cup\cdots\cup L_m\subset S^3$ be a smooth oriented link with zero pairwise linking numbers: $\operatorname{lk}(L_i,L_j)=0$ for $i\ne j$. Linking number is the signed intersection with a Seifert surface for the other component. For each component choose its zero-framed longitude $\lambda_i$ and a meridian $\mu_i$ in the boundary diagram, placed so that $\operatorname{lk}(\lambda_i,\mu_i)=0$, as in the source construction. Both auxiliary curves carry zero framing.

A four-dimensional round one-handle is $[-1,1]\times S^1\times D^2$, attached along the two solid tori $\{-1,1\}\times S^1\times D^2$. Let $R(L)$ be obtained from $B^4$ by attaching one such handle at the neighborhoods of $\lambda_i$ and $\mu_i$, with their specified framings, for each $i$. The original link components remain as curves in the new boundary. Call $L$ round-handle slice if its components bound pairwise disjoint proper locally flat disks in $R(L)$; locally flat means the disk embedding has the standard local topological plane or half-plane model.
\paragraph{Statement recorded in the supplied source.}
Is every link $L\subset S^3$ with zero pairwise linking numbers round-handle slice? Equivalently, must there exist a disjoint family of locally flat disks $D_1,\ldots,D_m\subset R(L)$ with $\partial D_i=L_i$, or is there a zero-pairwise-linking counterexample? Smooth slicing disks are not required. \sref{2923}{2}
\subsection{Short English statement}
Do all links with zero pairwise linking numbers become topologically slice after the specified round-one-handle attachments?
\subsection{Sources}
\begin{description}
\item[\hypertarget{source:2923:1}{S1}] ULAM.AI and the UnsolvedMath Contributors, UnsolvedMath: A Curated Collection of Open Mathematics Problems, record 2923 (KP-4.47). CC BY 4.0. \url{https://huggingface.co/datasets/ulamai/UnsolvedMath}.
\item[\hypertarget{source:2923:2}{S2}] R. İnanç Baykur, Robion C. Kirby and Daniel Ruberman, K3—A New Problem List in Low-Dimensional Topology, Mathematical Surveys and Monographs 295 (2026), author version, Problem 4.47 and Remarks, PDF page 228. \url{https://math.berkeley.edu/sites/default/files/surv-295-ruberman-watermarked-author-pdf.pdf}.
\end{description}
\label{end:2923}
\end{document}
